\section{Jaggedness de RL y las limitaciones fundamentales del modelo}

\subsection{¿Por qué el chiste no cambia?}

Karpathy plantea una observación intuitiva: «el chiste que contaba ChatGPT hace unos años sigue siendo el mismo chiste hoy» (Why do scientists not trust atoms? Because they make everything up.). El modelo avanza a pasos agigantados en las agentic task, pero se estanca en un dominio no optimizado por RL como «contar chistes».

\begin{importantbox}{On Rails vs Off Rails}
O bien te encuentras en un dominio en el que el modelo fue entrenado y optimizado (on rails) y disfrutas de la aceleración de una superinteligencia, o bien te encuentras en un dominio no optimizado (off rails) y experimentas la lentitud y la repetición del modelo. Esta dualidad es una característica fundamental del entrenamiento con RL: los labs solo pueden mejorar dominios con recompensas verificables.
\end{importantbox}

\subsection{La ausencia de generalización}

Karpathy no considera que la hipótesis de que «volverse mejor en código $\rightarrow$ volverse mejor en todo» se cumpla del todo. La mejora del modelo se concentra fuertemente en los dominios optimizados, y el avance en las soft skills es mínimo.

\subsection{Speciation: la especialización del modelo}

\begin{knowledgebox}{La analogía del reino animal}
Los cerebros de la naturaleza son extremadamente diversos: distintos animales tienen cortezas visuales hiperdesarrolladas u otras regiones especializadas. ¿Debería el modelo de IA también speciate (especializarse)? Karpathy considera que deberíamos ver más modelos especializados (como un modelo enfocado en matemáticas con Lean), pero por ahora los labs siguen persiguiendo el monoculture de «un solo modelo que hace todo». Parte de la razón es que la ciencia del fine-tuning sin pérdida de capacidades aún no es madura.
\end{knowledgebox}

\subsection{Resumen de la sección}

La jaggedness del modelo expone las limitaciones fundamentales del entrenamiento con RL: solo los dominios verificables se optimizan de manera efectiva. En el futuro tal vez se necesite más especialización de modelos (speciation), pero la inmadurez de la ciencia del fine-tuning y la preferencia de los labs por modelos de propósito general limitan esta tendencia.


